\documentclass{article}
\usepackage{graphicx} % Required for inserting images
\usepackage{amsmath}
\title{STT-2902}
\author{Julien Miron}


\begin{document}

\section*{Équipe 4 – Inférence - Intervalles de confiance}

Objectif : Expliquer le principe d’inférence statistique à partir d’un ou plusieurs échantillons dans le contexte des intervalles de confiance.


\vspace{1cm}
Suggestions :
\begin{itemize}
    \item Proposez une activité où les autres doivent estimer une valeur inconnue à partir d’un échantillon (ex. : moyenne d’une population de jetons numérotés).
    \item Ajoutez des intervalles de confiance, expliquez leur signification.

\end{itemize}


\vspace{1cm}
\textbf{Requis} :
\begin{enumerate}
    \item Vous devez clairement expliquer ce qu'un intervalle de confiance de niveau $(1-\alpha)$ nous donne comme information:
    
    $\rightarrow P( \text{le paramètre d'intérêt est dans l'intervalle de confiance} )=1-\alpha$. 
    \item Fournir un exemple visuel (graphique ou schéma) illustrant plusieurs intervalles de confiance pour différents échantillons.
\end{enumerate}

\end{document}
